The proof of Theorem~\ref{theorem_S_equal_T} will
be given
at the end of Section~\ref{section_permute}.

\section{The base case}
\label{section_base_case}

Let $\overline{\lambda}$ denote the
reverse of $\lambda$, that is,
$\overline{\lambda} = (\lambda_{n}, \ldots, \lambda_{2}, \lambda_{1})$.
\begin{lemma}
The sum $S(\lambda)$ can be expressed as
$$
\sum_{\sigma \in \Sigma(\lambda)}
(-1)^{|\sigma|} \cdot (-1)^{g(\sigma)} 
=
(-1)^{\binom{n}{2}}
\cdot
\sum_{\sigma \in \Sigma(\overline{\lambda})}
(-1)^{|\sigma|} \cdot (-1)^{f(\sigma)} . 
$$
\label{lemma_reverse}
\end{lemma}
\begin{proof}
Let $\tau_{0}$ denote the longest permutation
in $\Sym_{n}$,
that is,
$\tau_{0} = n \cdots 21$.
Note that $(-1)^{\tau_{0}} = (-1)^{\binom{n}{2}}$
and $\overline{\lambda} = \tau_{0} \lambda$.
For any ordered partition $\sigma \in \Qn$ 
we have the
following equality between permutations:
$\tau_{0} g(\sigma) = f(\tau_{0} \sigma)$.
In particular, their signs agree, that is,
$(-1)^{\binom{n}{2}} \cdot (-1)^{g(\sigma)}
= 
(-1)^{f(\tau_{0} \sigma)}$.
Now apply Lemma~\ref{lemma_reordering}
using the permutation $\tau_{0}$.
The result follows after summing over all faces
$\sigma \in \Sigma(\lambda)$.
\end{proof}


\begin{prop}
Let $\lambda \in \Rrr^{n}$ be a sequence such that
$\lambdaweaklyincreasing$.
Then the sum $S(\lambda)$ is given by
$$
S(\lambda)
=
\begin{cases}
(-1)^{n} & \text{ if } \lambda_{1} > 0, \\
0 & \text{ otherwise.}
\end{cases}
$$ 
\label{proposition_S_lambda_increasing}
\end{prop}
\begin{proof}
Lemma~\ref{lemma_reverse}
asserts that
$$
S(\lambda)
=
(-1)^{\binom{n}{2}}
\cdot
\sum_{\sigma\in\Sigma(\overline{\lambda})}
(-1)^{|\sigma|} \cdot (-1)^{f(\sigma)}.
$$
Moreover,
the entries of 
$\overline{\lambda}$
are in weakly decreasing order.
Hence we can use
the decomposition
in Proposition~\ref{proposition_Sigma_lambda_decomposition}
for the complex 
$\Sigma(\overline{\lambda})$,
that is, we have
$$
\Sigma(\overline{\lambda})
=
\bigcup_{\tau \in \mathcal{A}(\overline{\lambda})} [R(\tau), \tau] .
$$
Observe the function
$(-1)^{f(\sigma)}$
is constant on the interval
$[R(\tau), \tau]$.
Hence the sum
$$ \sum_{\sigma \in [R(\tau), \tau]} (-1)^{|\sigma|} \cdot (-1)^{f(\sigma)} $$
is zero unless $R(\tau) = \tau$.
But we can only have $R(\tau)=\tau$
when $\tau$ is the longest permutation
$n \cdots 21$. For this permutation
the sum is $(-1)^{n} \cdot (-1)^{\binom{n}{2}}$.
This permutation occurs in $\mathcal{A}(\overline{\lambda})$
if and only if $\lambda_{1} > 0$.
Hence the result follows.
\end{proof}

\begin{lemma}
The following two evaluations of $T(\lambda)$ hold:
\begin{enumerate}
\item\label{lemma7.3_1} If $\lambda_{1},\lambda_{2},\ldots,\lambda_{n}>0$, then $T(\lambda)=1$.
\item\label{lemma7.3_2} If $\lambda_{1}\leq 0$, then $T(\lambda)=0$.
\end{enumerate}
\label{lemma_T_two_statements}
\end{lemma}
\begin{proof}
We begin by proving the first statement.
By the hypothesis on $\lambda$, for every maximal matching~$p$,
we have $c(p,\lambda)=1$.
We prove 
$T(\lambda) = 1$ for all even $n$ by induction.
The induction basis $n=2$ is straightforward to check
since there is only one maximal matching on two vertices.
Assume now it is true for $n-2$.
Observe that the parity of the number 
of crossings to which the edge $\{1,j\}$ contributes is
the same as the parity of~$j$. 
Let $q \in M_{n-2}$ be the matching obtained from $p$ by removing
the vertices $1$ and $j$, and then applying 
the unique order-preserving relabeling of the remaining vertices.
This gives a bijection from the set of maximal matchings $p\in M_{n}$ such that
$1$ and $j$ are matched and the set of maximal matchings $q\in M_{n-2}$, and we
have $(-1)^p=(-1)^j (-1)^q$.
Hence the sum is given by
$$
T(\lambda)
=
\sum_{p \in M_{n}} (-1)^{p}
=
\sum_{j=2}^{n}  (-1)^{j} \cdot \sum_{q \in M_{n-2}} (-1)^{q}
=
\sum_{j=2}^{n}  (-1)^{j} 
=
1 ,
$$
where the third step is the induction hypothesis.

Assume now that $n$ is odd.
By using the sign-preserving bijection in
Remark~\ref{remark_add_extra_edge}
between $M_{n}$ and~$M_{n+1}$,
we obtain that
the two sums
$\sum_{p \in M_{n}} (-1)^{p}$
and
$\sum_{r \in M_{n+1}} (-1)^{r}$
are equal, and we have already seen that the second sum is equal to $1$.



We next prove the second statement.
Suppose that $\lambda_{1} \leq 0$. 
Then for every maximal matching~$p$, we claim $c(p,\lambda)=0$. 
Indeed, if the vertex $1$ is isolated then
$c_{1}(\lambda_{1}) = 0$ is a factor of $c(p, \lambda)$.
If not, let $e = \{1,j\}$ be the edge of $p$ containing the vertex $1$.
The contribution of $e$ to $c(p,\lambda)$ is $c_{2}(\lambda_{1},\lambda_{j})=0$
which is independent of the value of $\lambda_{j}$.
In both cases, $c(p,\lambda)$ has a factor equal to~$0$, so it is~$0$. 
\end{proof}

By combining
Proposition~\ref{proposition_S_lambda_increasing}
and
Lemma~\ref{lemma_T_two_statements},
we obtain the following corollary.
\begin{coro} 
If $\lambda_{1} \leq \lambda_{2} \leq \cdots \leq \lambda_{n}$, then
$S(\lambda)=(-1)^{n} \cdot  T(\lambda)$.
\label{corollary_base_case}
\end{coro}







\section{Permuting the entries of \texorpdfstring{$\lambda$}{lambda}}
\label{section_permute}

We now give a way to reduce the general case to the case treated
in the previous section,
assuming that we know the identity for smaller values of $n$.
This reduction is
directly inspired by Herb's paper~\cite{Herb} 
on discrete series characters. 

Let $\lambda \in \Rrr^{n}$
with no ordering hypothesis on $\lambda$.
Suppose that $n \geq 2$, and fix $1 \leq i \leq n-1$.
Recall that $s_{i}$ is the simple transposition $(i,i+1)$
in the symmetric group $\Sym_{n}$.
Also, we write
$\mu=(\lambda_{1},\ldots,\lambda_{i-1},\lambda_{i+2},\ldots,\lambda_{n})\in\R^{n-2}$.

\begin{prop}
The following two identities hold:
\begin{align}
S(\lambda) + S(s_{i}\lambda)
& =
-2 \cdot \ungras_{\lambda_{i}+\lambda_{i+1}>0}\cdot S(\mu) ,
\label{equation_S} \\
T(\lambda) + T(s_{i}\lambda)
& =
2 \cdot \ungras_{\lambda_{i}+\lambda_{i+1}>0}\cdot T(\mu) .
\label{equation_T}
\end{align}
\label{proposition_S_equal_T}
\end{prop}
\begin{proof}
We begin by proving~\eqref{equation_S}.
For $\sigma \in \Sigma_{n}$, let $s_{i}\sigma$ denote the
ordered partition where we exchange the elements $i$ and $i+1$.
Note that $s_{i}$ is an involution on $\Sigma_{n}$
and that it preserves the number of blocks.
We write
$\Sigma(\lambda) =\Sigma^{\prime}(\lambda) \sqcup \Sigma^{\prime\prime}(\lambda)$,
where $\Sigma^{\prime\prime}(\lambda)$ is the set of fixed points of $s_{i}$
in~$\Sigma(\lambda)$,
that is, 
$\Sigma^{\prime}(\lambda)$ is the set of $\sigma\in\Sigma(\lambda)$
such that $i$ and $i+1$ are in different blocks of $\sigma$
and
$\Sigma^{\prime\prime}(\lambda)$ is the set of $\sigma\in\Sigma(\lambda)$
such that $i$ and $i+1$ are in the same block of $\sigma$.



Note that the action of $s_{i}$ gives a bijection between
$\Sigma^{\prime}(\lambda)$ and $\Sigma^{\prime}(s_{i} \lambda)$.
Furthermore, for an ordered partition $\sigma$ in $\Sigma^{\prime}(\lambda)$
we have
$g(s_{i}\sigma) = s_{i}\cdot g(\sigma)$.
We obtain
\begin{align}
\sum_{\sigma \in \Sigma^{\prime}(\lambda)}
(-1)^{|\sigma|} \cdot (-1)^{g(\sigma)}
+
\sum_{\sigma \in \Sigma^{\prime}(s_{i}(\lambda))}
(-1)^{|\sigma|} \cdot (-1)^{g(\sigma)}
& = 0 
.
\label{equation_Sigma_prime}
\end{align}
We next consider $\Sigma^{\prime\prime}(\lambda)$.
Note that $\Sigma^{\prime\prime}(\lambda) =\Sigma^{\prime\prime}(s_{i}\lambda)$, so 
it remains to show that
\begin{align}
\sum_{\sigma\in\Sigma^{\prime\prime}(\lambda)}
(-1)^{|\sigma|} \cdot (-1)^{g(\sigma)}
& =
-\ungras_{\lambda_{i}+\lambda_{i+1}>0}\cdot S(\mu).
\label{equation_Sigma_prime_prime}
\end{align}


Suppose first that $\lambda_{i}+\lambda_{i+1}\leq 0$.
Define a matching on $\Sigma^{\prime\prime}(\lambda)$ in the following manner.
If $\sigma=(C_{1},C_{2},\ldots,C_{r}) \in \Sigma^{\prime\prime}(\lambda)$
satisfies
$i, i+1 \in C_{s}$ and $|C_{s}|\geq 3$, match it with
$$ \sigma^{\prime}
=
(C_{1},\ldots,C_{s-1},C_{s}-\{i,i+1\},\{i,i+1\},C_{s+1},\ldots,C_{r}) . $$
The hypothesis that $\lambda_{i}+\lambda_{i+1} \leq 0$
and Lemma~\ref{lemma_split_block}
imply that 
$\sigma^{\prime} \in \Sigma^{\prime\prime}(\lambda)$.
Furthermore, this is a perfect matching
since no ordered partition 
in $\Sigma^{\prime\prime}(\lambda)$
can begin with the block~$\{i,i+1\}$.
It is easy to see that 
$(-1)^{g(\sigma)}=(-1)^{g(\sigma^{\prime})}$ if $\sigma$ and $\sigma^{\prime}$ are matched.
So we obtain
\begin{align*}
\sum_{\sigma\in\Sigma^{\prime\prime}(\lambda)}
(-1)^{|\sigma|} \cdot (-1)^{g(\sigma)}
&=0.
%\label{equation_S_prime}
\end{align*}


Now assume that $\lambda_{i}+\lambda_{i+1}>0$.
We consider another matching on $\Sigma^{\prime\prime}(\lambda)$, defined as follows.
If $\sigma=(C_{1}, C_{2}, \ldots, C_{r})\in\Sigma^{\prime\prime}(\lambda)$
satisfies $i, i+1 \in C_{s}$ and $|C_{s}| \geq 3$, match it with
$$ \sigma^{\prime}
= 
(C_{1},\ldots,C_{s-1},\{i,i+1\},C_{s}-\{i,i+1\},C_{s+1},\ldots,C_{r}). $$
The hypothesis that $\lambda_{i}+\lambda_{i+1} > 0$
and Lemma~\ref{lemma_split_block}
imply that 
$\sigma^{\prime} \in \Sigma^{\prime\prime}(\lambda)$.
The unmatched elements
are the ordered partitions whose last block is $\{i, i+1\}$.
Again, it is straightforward to see that 
$(-1)^{g(\sigma)}=(-1)^{g(\sigma^{\prime})}$ if $\sigma$ and $\sigma^{\prime}$ are matched.
Denote by $\Sigma^{\prime\prime\prime}(\lambda)$ the set of unmatched elements
in~$\Sigma^{\prime\prime}(\lambda)$.
We obtain
$$ \sum_{\sigma\in\Sigma^{\prime\prime}(\lambda)}
(-1)^{|\sigma|} \cdot (-1)^{g(\sigma)}
=
\sum_{\sigma\in\Sigma^{\prime\prime\prime}(\lambda)}
(-1)^{|\sigma|} \cdot (-1)^{g(\sigma)} . 
$$
Identify
$[n-2]$ and $[n]-\{i,i+1\}$
using the unique order-preserving bijection between these sets.
We note that
$\sigma=(C_{1}, \ldots, C_{r},\{i,i+1\})\in\Sigma_{n}$
is an element of $\Sigma(\lambda)$ if and only if
$\tau := (C_{1},\ldots,C_{r})$ is an element of $\Sigma(\mu)$.
This induces a bijection $\Sigma^{\prime\prime\prime}(\lambda)\iso\Sigma(\mu)$.
Also, we have $|\sigma|=|\tau|+1$ and $(-1)^{g(\sigma)}=(-1)^{g(\tau)}$.
So finally we find that
\begin{align*}
\sum_{\sigma\in\Sigma^{\prime\prime\prime}(\lambda)}
(-1)^{|\sigma|} \cdot (-1)^{g(\sigma)}
&=
-
\sum_{\tau\in\Sigma(\mu)}
(-1)^{|\tau|} \cdot (-1)^{g(\tau)}
=
-S(\mu).
%\label{equation_S_prime_prime}
\end{align*}
Identity~\eqref{equation_S}
follows by combining equation~\eqref{equation_Sigma_prime} 
with twice equation~\eqref{equation_Sigma_prime_prime}
in the two cases
$\Sigma^{\prime\prime}(\lambda)$ and $\Sigma^{\prime\prime}(s_{i} \lambda)$.


Identity~\eqref{equation_T} 
is proved as in Case~I of Lemma~2.14 in~\cite{Herb}.
Let $M_{n}^{\prime\prime}$ be the set
of maximal matchings in $M_{n}$ where
$\{i,i+1\}$ is an edge and
let $M_{n}^{\prime}$ be the complement,
that is, where the vertices $i$ and $i+1$ are not matched with each other.

For $p \in M_{n}$, let $s_{i}p$ denote the
matching where we exchange the vertices $i$ and~$i+1$.
Note that $s_{i}$ is an involution on the
set $M_{n}$
and that $M_{n}^{\prime\prime}$ is the set of fixed points
of~$s_{i}$.
Furthermore, for a matching $p \in M_{n}^{\prime}$
we have
$(-1)^{s_{i} p} = - (-1)^{p}$
and
$c(s_{i}p, s_{i} \lambda) = c(p, \lambda)$.
Hence we obtain
\begin{align}
\sum_{p \in M_{n}^{\prime}} (-1)^{p} \cdot c(p, \lambda)
+
\sum_{p \in M_{n}^{\prime}} (-1)^{p} \cdot c(p, s_{i} \lambda)
= 0 .
\label{equation_M_prime}
\end{align}
Again we use
the unique order-preserving bijection $[n-2] \iso [n] - \{i,i+1\}$.
Thus if $p$ is a matching in $M_{n}^{\prime\prime}$, we can view
$p-\{i,i+1\}$ as a matching
on the vertex set $[n-2]$, which we will denote by~$q$.
This gives a bijection between
$M_{n}^{\prime\prime}$
and $M_{n-2}$.
Note that the two matchings $p$ and $q$ have the same sign,
that is, $(-1)^{p} = (-1)^{q}$.
Furthermore
\begin{align*}
c(p,\lambda) + c(p,s_{i}\lambda) 
& =
(c_{2}(\lambda_{i},\lambda_{i+1}) + c_{2}(\lambda_{i+1},\lambda_{i}))
\cdot
c(q,\mu) \\
& =
2 \cdot \ungras_{\lambda_{i}+\lambda_{i+1} > 0}
\cdot
c(q,\mu) .
\end{align*}
Now summing over all $p \in M_{n}^{\prime\prime}$
we obtain
\begin{multline}
\label{equation_M_prime_prime}
\sum_{p \in M_{n}^{\prime\prime}}
(-1)^{p} \cdot c(p, \lambda)
+
\sum_{p \in M_{n}^{\prime\prime}}
(-1)^{p} \cdot c(p, s_{i} \lambda)
\\ = 
2 \cdot \ungras_{\lambda_{i}+\lambda_{i+1} > 0} 
\cdot
\sum_{q \in M_{n-2}} (-1)^{q} \cdot c(q, \mu) .
%\nonumber
\end{multline}
Identity~\eqref{equation_T} 
now follows by summing equations~\eqref{equation_M_prime}
and~\eqref{equation_M_prime_prime}.
\end{proof}


We now include the proof of Theorem~\ref{theorem_S_equal_T}. 

